\chapter{Discussion}
\label{chap:discussion}

\section{Purpose of the Discussion}
The results chapters report what was measured. This chapter interprets those
measurements against the research questions, the design objectives, and the
claim boundaries. The current v0.0.7 run adds a larger diagnostic population,
a paired role-contract correction, and a numerical benchmark. Together these
show that the artifact can preserve source evidence,
provider provenance, semantic uncertainty, and execution health across a
large frozen population. Those properties make later human validation
possible; they do not replace it.

\section{Synthesis of the Four Research Questions}
\subsection{RQ1: Broad Coverage Requires Explicit Ownership}
The 76-concept catalogue is broader than any one provider because its purpose
is to operationalise a review checklist, not reproduce a linter. Native
AnsibleGuard detectors cover repository structure, role relationships,
runtime-variable evidence, testing/documentation signals, and policy-oriented
checks. ansible-lint remains authoritative for established Ansible idioms such
as FQCN and naming conventions. GitLeaks contributes candidate-secret pattern
evidence. Where two providers describe the same underlying concern,
canonicalisation keeps one review occurrence while retaining both evidence
records.

This ownership model avoids two opposite errors. Reimplementing all
ansible-lint rules would create unnecessary divergence from a mature tool;
delegating the complete checklist would leave repository semantics and review
workflow concerns uncovered. The 15 Yes, 18 Partial, and 5 No checklist mapping
also prevents the catalogue size from being mistaken for complete automation.
Human concerns such as architectural appropriateness remain human concerns.

\subsection{RQ2: Completion Is Not the Same as Provider Success}
All 149 repository analyses completed, yet ansible-lint was partial on 21
repositories and the run retained 200 diagnostics. This is a stronger result
than reporting 149 successful process exits. It demonstrates that repository
orchestration survived provider-local problems and that the output tells a
consumer where evidence is incomplete. A provider timeout, missing dependency,
configuration failure, or YAML parse problem is an environment/execution fact;
it is not an ordinary source smell and must not contribute to a quality score.

The separation has direct consequences for automation. A CI policy may fail on
an error-level occurrence, retry on a provider tool error, warn on partial
coverage, or accept a documented unavailable optional provider. Treating all
four cases as findings would make remediation misleading. Treating them all as
zero findings would make the scan falsely reassuring.

\subsection{RQ3: Version Deltas Reflect Evidence Semantics}
The paired comparison holds repositories, SHAs, selected files, and LOC
constant. Consequently, the fall from 331,123 to 328,385 raw findings and the
rise from 297,557 to 307,951 canonical occurrences cannot be explained by a
different corpus. They arise from changed emission, identity, provenance, and
semantic classification behaviour.

In particular, v0.0.6 avoids unsafe matching of generic location-only findings
across files, preserves repeated assignment evidence where appropriate, and
names dynamic/runtime limitations as \code{UNCERTAIN}. The increase from
19,431 to 43,669 uncertain occurrences is therefore not evidence that source
quality suddenly worsened. Conversely, the decrease in non-uncertain evidence
from 278,126 to 264,282 is not proof that precision improved. Both are
observable effects of a more conservative evidence model; accuracy still
requires labelled examples.

\subsection{RQ4: The Variable Problem Is Now Mechanism-Specific}
The variable detector reports 43,035 occurrences, of which 41,351 are
\code{UNCERTAIN} and 1,684 are \code{UNRESOLVED}. This split changes the
research question from ``why is the total large?'' to ``which evidence source
is absent or dynamic in each mechanism?'' Static include/import propagation,
lexically scoped Jinja loop variables, ordered persistent definitions, and
guard modelling remove demonstrated semantic errors without leaking symbols
to disconnected files.

The remaining unresolved set is concentrated. SAP, Linuxfabrik, Infrared,
Forklift, and CMS account for 1,314 of 1,684 cases (78.03\%). Their dominant
mechanisms---dynamic \code{include_vars}, role-entry contracts, composed
variables, plugin inputs, extra variables, and external derivations---suggest
targeted improvements. A global suppression would hide useful strict cases;
mechanism-aware producers or uncertainty reasons are the defensible next step.

\section{Interpreting the Main Result Distributions}
\subsection{Why Style Dominates}
Style contributes 127,420 canonical occurrences (41.38\%), and
\code{missing-tags} alone contributes 83,446. This dominance is partly a
property of granularity: a repetitive task-level convention can fire many
times, while a repository-level architecture concern may fire once. It does
not follow that style is the most harmful problem in the corpus.

For this reason, the complete \code{technical-full} profile is appropriate for
research inventory and deep auditing, but it should not be the only view shown
to every user. Security, reliability, quality, migration, or PR profiles let a
consumer align evidence with a decision. Counts should also be grouped by rule
and repository before prioritisation so a high-frequency convention does not
visually erase a low-frequency error.

\subsection{Severity, Status, and Diagnostics Are Independent Axes}
The 184,323 warnings, 116,729 notes, and 6,899 errors express remediation
priority assigned to source evidence. Semantic status expresses what the
detector knows about the predicate. Diagnostics express whether analysis
executed as intended. These axes answer different questions:
\begin{itemize}
  \item severity asks how strongly policy should react;
  \item status asks whether evidence is confirmed by the detector, uncertain,
  exposed, unresolved, or ordered in conflict;
  \item diagnostic level asks how analysis health should be communicated.
\end{itemize}
An error-severity occurrence can still be uncertain. A warning can be
confirmed. A provider tool error is neither of those source judgements. Output
consumers must preserve this multidimensional model rather than collapse it to
one score.

\subsection{Canonicalisation Is Evidence Compression, Not Deletion}
The difference between 328,385 raw findings and 307,951 canonical occurrences
is 20,434 records, or 6.22\%. These records were not declared false positives.
They were merged when multiple pieces of provider evidence supported one
review occurrence. The contributing evidence remains available for audit.
Canonicalisation therefore reduces review duplication while allowing the user
to inspect which provider, rule, coordinate, message, and fingerprint produced
the conclusion.

\section{Evidence-to-Decision Traceability}
Table~\ref{tab:evidence_decision} translates major observations into defensible
actions. This is intentionally more conservative than ranking repositories by
their total count.

\begin{longtable}{p{4.0cm}p{5.1cm}p{5.1cm}}
\caption{From observed evidence to interpretation and action.}
\label{tab:evidence_decision}\\
\toprule Observation & Defensible interpretation & Recommended action \\
\midrule\endfirsthead
\toprule Observation & Defensible interpretation & Recommended action \\
\midrule\endhead
High total dominated by \code{missing-tags} & Repeated policy/style convention,
not an equal number of severe defects & Use a profile, group by rule, and
configure tag policy before remediation. \\
\code{UNRESOLVED} variable & No modelled definition, guard, producer, or named
uncertainty source is visible & Inspect source context and role/inventory
contract; fix source or add demonstrated producer semantics. \\
\code{UNCERTAIN} variable & A concrete dynamic or external limitation prevents
a strict resolution claim & Review the recorded uncertainty reason; do not
label it false positive automatically. \\
ansible-lint \code{partial} & Some delegated evidence is incomplete & Inspect
provider diagnostics, dependency/configuration, and root plan; do not report a
clean delegated result. \\
GitLeaks candidate & A secret-like pattern is present & Validate deployability,
test/example context, and rotation need; preserve original evidence. \\
Raw count falls while canonical rises & Provider emission and occurrence
identity changed differently & Compare rule/provider/status deltas on the
frozen population; do not infer quality improvement from either total alone. \\
One repository dominates a rule & Mechanism or project convention is
concentrated & Sample that repository separately and avoid generalising its
pattern to the ecosystem. \\
\bottomrule
\end{longtable}

\section{Implications for Stakeholders}
\textbf{Developers} need short, local feedback. They should begin with a fast
or quality profile and inspect exact source context, then escalate uncertain
semantic cases rather than suppress the entire rule. \textbf{Pull-request
reviewers} need introduced/fixed/unchanged classification, with file-aware
identity and review-range rules. \textbf{Maintainers} need complete repository
structure, dependency, documentation, and testing evidence.

\textbf{Security engineers} need security-focused findings plus GitLeaks
provenance, but must distinguish candidate patterns from confirmed exploitable
secrets. \textbf{CI/platform teams} need stable JSON/SARIF, schema validation,
exit policy, and provider health. \textbf{Researchers and supervisors} need
immutable manifests, denominators, version and configuration identities, aggregate
CSV, per-repository envelopes, and explicit limitations. The same analyzer can
serve these users only because profiles change the view while the underlying
evidence contract remains stable.

\section{Implications for Static-Analysis Design}
Three broader design conclusions follow. First, infrastructure-as-code analysis
needs an artifact model: a defaults file, task file, inventory, and CI fixture
do not have interchangeable semantics. Second, cross-file reasoning should be
graph-constrained. Repository-wide symbol unions reduce apparent findings but
introduce unsound leakage. Third, uncertainty should be first-class data. A
detector that cannot resolve a dynamic include should name that limitation,
not force a binary clean/defect answer.

Provider orchestration also belongs to the analysis model. The location and
configuration of ansible-lint project roots affect both duplication and
coverage. Collapsing overlapping unconfigured roots while preserving
independently configured roots is therefore correctness work, not only an
execution optimisation. Similarly, packaging the canonical schema and copying
it into outputs is part of reproducibility: a wheel installed outside the
source checkout must validate exactly the same contract.

\section{Threats to Validity and Mitigations}
\begin{longtable}{p{3.0cm}p{5.4cm}p{5.7cm}}
\caption{Validity threats, current mitigation, and residual risk.}
\label{tab:validity_matrix}\\
\toprule Type & Threat & Mitigation and residual limitation \\
\midrule\endfirsthead
\toprule Type & Threat & Mitigation and residual limitation \\
\midrule\endhead
Construct & A finding count may be mistaken for defects or quality. & Terms are
defined as detector evidence; status, severity, diagnostics, and canonical
units are separate. Human correctness remains unmeasured. \\
Internal & Version differences may result from population drift. & The paired
comparison fixes 149 SHAs, 27,709 files, and 1,178,908 LOC. Tool/config changes
must still be recorded when reproducing. \\
Internal & Provider failure may masquerade as absence of findings. & Health and
diagnostics are explicit and repository execution is isolated. A partial run
still lacks some delegated evidence. \\
External & Mature public repositories may not represent all Ansible use. & The
eligibility pipeline and attrition evidence are documented. Results are not
claimed as ecosystem prevalence. \\
Conclusion & Large totals may create spurious accuracy claims. & Descriptive
statistics, mechanism samples, and claim boundaries are used. No precision or
recall is reported without an adjudicated benchmark. \\
Reliability & Local checkout assumptions may break installed execution. & The
schema is packaged and wheel-install gates run outside the checkout. External
provider availability remains environment-dependent. \\
Researcher bias & Mechanism labels may reflect one analyst. & Samples retain
source/provenance for audit; the proposed next evaluation uses two blind
reviewers, adjudication, and agreement reporting. \\
\bottomrule
\end{longtable}

\section{How the Versioned Evidence Fits Together}
The v0.0.6 chapter preserves the complete 149-repository experiment and its
paired v0.0.5 comparator. The current v0.0.7 chapter uses a 207-repository
frozen manifest, but its production aggregate has 206 analyzable repositories.
The paired role-contract run holds those 206 outputs fixed. The corrected
threshold benchmark separately covers all 207 repositories. These populations
answer different questions: a change in total between v0.0.6 and v0.0.7
cannot be attributed solely to a detector version, while the paired 206-run
comparison can isolate the role-contract mechanism. The benchmark's 207-of-207
completion describes metric extraction, not provider execution health.

The single-agent variable review narrows plausible error mechanisms but is not
an independent accuracy benchmark. For the numerical thresholds, the 24
frozen agreement pilots and weighted primary packets are ready; no double
human labels have been entered. In particular, the Q3 large-task candidate
would more than double the current review workload, but no quality estimate
yet justifies adopting it.

\section{Priorities for the Next Iteration}
The next iteration should be governed by labelled evidence and concentrated
mechanisms rather than by reducing the largest number. The order is:
\begin{enumerate}
  \item create a reproducible stratified sample across concept, provider,
  status, severity, and repository concentration;
  \item define TP, FP, and borderline guidance per sampled rule, use two
  independent reviewers, adjudicate, and report agreement;
  \item refine demonstrated variable producers such as inventory/plugin inputs
  and dynamic \code{include_vars} families, then investigate the Miarka miss;
  \item build a mutation or curated benchmark for recall-sensitive rules;
  \item keep provider health, schema portability, remote equivalence, and the
  frozen population as non-regression gates;
  \item publish exact versions, configurations, manifests, schema hashes, and
  labelled evidence with the next result set.
\end{enumerate}

This priority order avoids optimising for a cosmetically smaller count. A rule
should change only when source evidence demonstrates an incorrect predicate,
missing semantic mechanism, unsuitable default policy, or output-contract
problem.
